% Unit 044 terjemahan Bahasa Indonesia dari chapter3.tex baris 2553--2713.
% Sumber: Methods of Algebra, Volume 2, commit
% 9a5803ff77dd3257484cb177f851a73770a59dd3 (CC BY 4.0).
% stable-unit-id: o014.aljabr2.chapter3.example-lim1
% Cakupan: Bagian 3.13 lengkap.

% segment-id: o014.li2.chapter3.example-lim1.h001
\section{Contoh: \texorpdfstring{$\lim\nolimits^1$}{lim1}}\label{sec:lim1}
\hypertarget{o014.aljabr2.chapter3.example-lim1}{ }

% segment-id: o014.li2.chapter3.example-lim1.p001
Sebagai contoh pertama funktor turunan, kita pilih $\varprojlim$ yang berasal
dari suatu kelas yang umum dijumpai. Pandang $\Z_{\geq 0}$ sebagai kategori
melalui struktur urutan total standarnya; dengan demikian, kategori
$\Z_{\geq 0}^{\opp}$ dapat ditampilkan sebagai
$\cdots \to 2 \to 1 \to 0$. Untuk sebarang kategori $\mathcal{A}$, tinjau
kategori funktor
% segment-id: o014.li2.chapter3.example-lim1.q001
\begin{equation}\label{eqn:InvSys}
	\index[sym1]{InvSys@$\cate{InvSys}$}
	\cate{InvSys}(\mathcal{A}) := \mathcal{A}^{\Z_{\geq 0}^{\opp}},
\end{equation}
yang objek-objeknya dapat diidentifikasi dengan data $(A_n, f_n)_{n \geq 0}$,
dengan $A_n \in \Obj(\mathcal{A})$ dan
$f_n \in \Hom(A_{n+1}, A_n)$; data seperti ini juga disebut
\emph{sistem invers} di $\mathcal{A}$. Bagian ini akan menggunakan funktor
turunan untuk mengkaji $\varprojlim$ dari sistem-sistem tersebut dalam kasus
kategori abelian.

% segment-id: o014.li2.chapter3.example-lim1.cn001
\begin{konvensi}\label{con:exact-product}
	\index{kategori abelian!mempunyai produk atau koproduk terhitung eksak (with exact countable products or coproducts)}
Jika kategori abelian $\mathcal{A}$ mempunyai produk terhitung (atau
koproduk terhitung, yakni jumlah langsung), dan funktor produk
$\prod: \mathcal{A}^{\Z_{\geq 0}} \to \mathcal{A}$ (atau funktor jumlah
langsung $\oplus: \mathcal{A}^{\Z_{\geq 0}} \to \mathcal{A}$) merupakan
funktor eksak, maka $\mathcal{A}$ dikatakan mempunyai \emph{produk
terhitung eksak} (atau \emph{koproduk terhitung eksak}).
\end{konvensi}

% segment-id: o014.li2.chapter3.example-lim1.p002
Funktor produk selalu eksak kiri. Oleh karena itu, dengan asumsi bahwa produk
terhitung ada, keeksakannya ekuivalen dengan sifat bahwa $\prod$
mempertahankan epimorfisme: untuk sebarang keluarga epimorfisme
$(g_n: A_n \to B_n)_{n \geq 0}$, morfisme yang bersesuaian
$\prod_{n \geq 0} g_n: \prod_{n \geq 0} A_n \to
\prod_{n \geq 0} B_n$ juga merupakan epimorfisme. Kasus koproduk terhitung
bersifat dual.

% segment-id: o014.li2.chapter3.example-lim1.hy001
\begin{hipotesis}\label{hyp:lim1}
Semua kategori abelian $\mathcal{A}$ yang ditinjau dalam bagian ini
diasumsikan mempunyai produk terhitung eksak.
\end{hipotesis}

% segment-id: o014.li2.chapter3.example-lim1.p003
Kategori modul $R\dcate{Mod}$ atas sebarang gelanggang $R$ mempunyai sifat
ini.

% segment-id: o014.li2.chapter3.example-lim1.de001
\begin{definisi}\label{def:Delta-diff}
	\index[sym1]{DeltaA@$\Delta_A$}
Diberikan objek $A = (A_n, f_n)_n$ dari
$\cate{InvSys}(\mathcal{A})$, definisikan morfisme pergeseran
$T_A: \prod_{n \geq 0} A_n \to \prod_{n \geq 0} A_n$ sebagai berikut:
jika $\mathcal{A} = \cate{Ab}$, maka
$T_A((a_n)_{n \geq 0}) := (f_n(a_{n+1}))_{n \geq 0}$; kasus umum
didefinisikan secara serupa. Definisikan pula morfisme kanonik
% segment-id: o014.li2.chapter3.example-lim1.q002
\[ \Delta_A := T_A - \identity: \prod_{n \geq 0} A_n \to \prod_{n \geq 0} A_n, \]
Ketika $A$ bervariasi, konstruksi ini memberikan endomorfisme $T$ dan
$\Delta = T - \identity$ dari funktor produk $\prod_{n \geq 0}$.
\end{definisi}

% segment-id: o014.li2.chapter3.example-lim1.p004
Perhatikan bahwa definisi di atas hanya memerlukan bahwa $\mathcal{A}$
merupakan kategori-$\cate{Ab}$ yang mempunyai produk terhitung.

% segment-id: o014.li2.chapter3.example-lim1.p005
Ingat bahwa $\varprojlim$ pada umumnya dibangun melalui penyama dan produk.
Dalam situasi sekarang, $\cate{InvSys}(\mathcal{A})$ secara alami merupakan
kategori abelian (Proposisi \ref{prop:functor-cat-Abel}), sehingga konstruksi
tersebut menjadi
$\varprojlim A = \Ker\left[ \Delta_A:
\prod_n A_n \to \prod_n A_n \right]$.

% segment-id: o014.li2.chapter3.example-lim1.p006
Selain itu, $\cate{InvSys}(\mathcal{A})$ juga mempunyai produk terhitung
eksak, sebab limit dalam kategori funktor dapat dibangun suku demi suku;
lihat \S\ref{sec:limit-functor}.

% segment-id: o014.li2.chapter3.example-lim1.le001
\begin{lema}
Konstruksi di atas memberikan funktor aditif eksak kiri
$\varprojlim: \cate{InvSys}(\mathcal{A}) \to \mathcal{A}$. Di sisi lain,
funktor produk
$\prod_{n \geq 0}: \cate{InvSys}(\mathcal{A}) \to \mathcal{A}$
merupakan funktor eksak.
\end{lema}
% segment-id: o014.li2.chapter3.example-lim1.pf001
\begin{bukti}
Aditivitas kedua funktor tersebut jelas, sedangkan keeksakan kiri
$\varprojlim$ mengikuti dari Contoh \ref{eg:limit-exactness}. Berdasarkan
hipotesis, funktor
$\prod_{n \geq 0}: (A_n, f_n)_n \mapsto \prod_n A_n$ juga
mempertahankan epimorfisme, sebab sifat itu berlaku pada setiap suku.
Karena itu, Catatan \ref{rem:exactness-mono-epi} menunjukkan bahwa funktor
tersebut eksak.
\end{bukti}

% segment-id: o014.li2.chapter3.example-lim1.de002
\begin{definisi}
	\index[sym1]{lim1@$\lim\nolimits^1$}
Definisikan keluarga funktor aditif
$\lim^n: \cate{InvSys}(\mathcal{A}) \to \mathcal{A}$ sebagai berikut:
% segment-id: o014.li2.chapter3.example-lim1.q003
\begin{align*}
	\lim\nolimits^0 & := \varprojlim = \Ker \Delta: \; A \mapsto \Ker \Delta_A, \\
	\lim\nolimits^1 & := \Coker \Delta: \; A \mapsto \Coker \Delta_A , \\
	\lim\nolimits^n & := 0, \quad n \in \Z_{\geq 2}.
\end{align*}
\end{definisi}

% segment-id: o014.li2.chapter3.example-lim1.p007
Sekarang misalkan $0 \to X \to Y \to Z \to 0$ merupakan barisan eksak
pendek di $\cate{InvSys}(\mathcal{A})$; dengan kata lain, barisan itu eksak
pada setiap suku. Terapkan funktor eksak $\prod_{n \geq 0}$ beserta
endomorfismenya $\Delta$ untuk memperoleh diagram komutatif dengan
baris-baris eksak
% segment-id: o014.li2.chapter3.example-lim1.g001
\[
\begin{tikzcd}
	0 \arrow[r] & \prod_n X_n \arrow[r] \arrow[d, "\Delta_X"'] & \prod_n Y_n \arrow[r] \arrow[d, "\Delta_Y"] & \prod_n Z_n \arrow[d, "\Delta_Z"] \arrow[r] & 0 \\
	0 \arrow[r] & \prod_n X_n \arrow[r] & \prod_n Y_n \arrow[r] & \prod_n Z_n \arrow[r] & 0 .
\end{tikzcd}
\]
Dengan menerapkan Teorema \ref{prop:snake-lemma}, diperoleh barisan eksak
panjang
% segment-id: o014.li2.chapter3.example-lim1.q004
\[ 0 \to \lim\nolimits^0 X \to \lim\nolimits^0 Y \to \lim\nolimits^0 Z \xrightarrow[\text{morfisme penghubung}]{\delta^0} \lim\nolimits^1 X \to \lim\nolimits^1 Y \to \lim\nolimits^1 Z \to 0 ; \]
barisan ini funktorial terhadap barisan eksak pendek. Untuk $n > 0$,
tetapkan pula $\delta^n := 0$. Hasilnya dapat dirangkum sebagai berikut.

% segment-id: o014.li2.chapter3.example-lim1.le002
\begin{lema}
Data $\left( \lim^n, \delta^n \right)_{n \geq 0}$ memberikan
funktor-$\delta$ kohomologis dari $\cate{InvSys}(\mathcal{A})$ ke
$\mathcal{A}$ (Definisi \ref{def:delta-functor}) yang memenuhi
$\lim^0 = \varprojlim$.
\end{lema}

% segment-id: o014.li2.chapter3.example-lim1.p008
Sekarang beralihlah ke funktor turunan kanan dari $\varprojlim$. Jika
$\mathcal{A}$ mempunyai cukup banyak objek injektif, maka
$\cate{InvSys}(\mathcal{A})$ juga demikian. Sesungguhnya, latihan-latihan
dalam Bab~\sourcecrossref{sec:Abel-cat}{2} telah membahas kasus umum kategori
funktor $\mathcal{A}^{\mathcal{C}}$, sedangkan yang diperlukan di sini ialah
kasus khusus $\mathcal{C} = \Z_{\geq 0}^{\opp}$. Untuk kasus ini, kita
langsung membuktikan hasil yang sedikit lebih kuat berikut.

% segment-id: o014.li2.chapter3.example-lim1.le003
\begin{lema}\label{prop:lim1-prep}
Dengan Hipotesis \ref{hyp:lim1}, definisikan subkategori penuh aditif
$\mathcal{R}$ dari $\cate{InvSys}(\mathcal{A})$ melalui
% segment-id: o014.li2.chapter3.example-lim1.q005
\[ \Obj(\mathcal{R}) = \left\{\begin{array}{l|l}
		R = (R_k, r_k)_{k \geq 0} & R \;\text{merupakan objek injektif}, \; \Delta_R \;\text{merupakan epimorfisme}, \\
		\; \in \Obj\left( \cate{InvSys}(\mathcal{A}) \right) & \forall k, \; R_k \;\text{merupakan objek injektif di $\mathcal{A}$}
\end{array}\right\}, \]
Maka, untuk setiap
$A \in \Obj\left(\cate{InvSys}(\mathcal{A})\right)$, terdapat
$R \in \Obj(\mathcal{R})$ beserta monomorfisme $A \hookrightarrow R$.
\end{lema}

% segment-id: o014.li2.chapter3.example-lim1.pf002
\begin{bukti}
Untuk setiap $m \in \Z_{\geq 0}$, definisikan funktor
$\mathcal{R}_m: \mathcal{A} \to \cate{InvSys}(\mathcal{A})$ yang
memetakan objek $X$ ke
% segment-id: o014.li2.chapter3.example-lim1.g002
\begin{equation}\label{eqn:Rm-construction}\begin{tikzcd}[row sep=tiny]
		\cdots \arrow[r, "{\identity}"] & X \arrow[r, "{\identity}"] & X \arrow[r] & 0 \arrow[r] & \cdots \arrow[r] & 0. \\
		\cdots & \text{suku ke-$m+1$} & \text{suku ke-$m$} & & &
\end{tikzcd}\end{equation}
Jelas bahwa funktor ini merupakan adjoin kanan dari
$\mathrm{ev}_m: (A_n, f_n)_{n \geq 0} \mapsto A_m$. Karena itu,
funktor tersebut membawa objek injektif ke objek injektif
(Proposisi \ref{prop:adjoint-injective-projective}). Selain itu, mudah
diperiksa bahwa $\Delta_{\mathcal{R}_m X}$ merupakan epimorfisme; suatu
invers kanannya bahkan dapat dituliskan secara eksplisit.

Diberikan $A = (A_n, f_n)_{n \geq 0}$, untuk setiap $m$ pilih objek
injektif $I_m$ dari $\mathcal{A}$ beserta monomorfisme
$A_m \hookrightarrow I_m$. Maka,
$A \hookrightarrow \prod_{m \geq 0} \mathcal{R}_m (I_m) =: R$.
Berdasarkan Lema \ref{prop:prod-injective-projective}, ruas kanan tetap
merupakan objek injektif. Bahkan, berdasarkan konstruksi, $R_k$
merupakan objek injektif dari $\mathcal{A}$ untuk setiap $k$. Karena
$\prod_m$ eksak,
$\Delta_R = \prod_m \Delta_{\mathcal{R}_m(I_m)}$ tetap merupakan
epimorfisme. Jadi, $R \in \Obj(\mathcal{R})$.
\end{bukti}

% segment-id: o014.li2.chapter3.example-lim1.p009
Dengan demikian, kita dapat membicarakan
$\mathrm{R}^n \varprojlim$ untuk $n \geq 0$. Ingat bahwa
$\mathrm{R}^n \varprojlim$ dapat dicirikan sebagai funktor-$\delta$
kohomologis universal (Korolari \ref{prop:derived-erasable}).

% segment-id: o014.li2.chapter3.example-lim1.th001
\begin{teorema}[S.\ Eilenberg]\label{prop:lim1}
Dengan Hipotesis \ref{hyp:lim1}, untuk $n > 0$ funktor
$\lim\nolimits^n$ dapat dihapus dalam pengertian Definisi
\ref{def:erasable}. Sebagai akibatnya, terdapat isomorfisme kanonik
$\mathrm{R}^1 \varprojlim \simeq \lim^1$, sedangkan
$\mathrm{R}^n \varprojlim = 0$ untuk $n \notin \{0, 1\}$.
\end{teorema}
% segment-id: o014.li2.chapter3.example-lim1.pf003
\begin{bukti}
Lema \ref{prop:lim1-prep} menunjukkan bahwa $\lim^1$ dapat dihapus.
Karena itu, $(\lim^n)_{n \geq 0}$ juga merupakan funktor-$\delta$
universal (Proposisi \ref{prop:erasable-univ-delta}).
\end{bukti}

% segment-id: o014.li2.chapter3.example-lim1.p010
Setelah $\lim^1$ diidentifikasi sebagai penghalang bagi keeksakan
$\varprojlim$, sekarang kita kaji keadaan-keadaan yang menjamin
$\lim^1 = 0$.

% segment-id: o014.li2.chapter3.example-lim1.ex001
\begin{contoh}\label{eg:quotient-lim1-0}
Jika objek $A$ dari $\cate{InvSys}(\mathcal{A})$ memenuhi
$\lim^1 A = 0$, maka setiap objek kuosien $Q$ dari $A$ juga memenuhi
$\lim^1 Q = 0$. Hal ini karena barisan eksak pendek
$0 \to K \to A \to Q \to 0$ menginduksi barisan eksak panjang yang
berakhir dengan $\lim^1 A \to \lim^1 Q \to 0$.
\end{contoh}

% segment-id: o014.li2.chapter3.example-lim1.ex002
\begin{contoh}\label{eg:lim1-section}
Jika setiap $f_n: A_{n+1} \to A_n$ mempunyai penampang (yakni invers
kanan) $s_n$, maka $\Delta_A$ juga mempunyai penampang
$\Sigma_A: \prod_n A_n \to \prod_n A_n$, sehingga $\Delta_A$
merupakan epimorfisme. Dalam notasi elemen, $\Sigma_A$ dapat dipilih
secara rekursif melalui
$\Sigma_A((b_n)_n) = (a_n)_n$, dengan $a_0 = 0$ dan
$a_{n+1} := s_n\left( a_n + b_n \right)$. Kasus umum mengikuti secara
serupa.
\end{contoh}

% segment-id: o014.li2.chapter3.example-lim1.de003
\begin{definisi}[Syarat Mittag-Leffler]\label{def:ML}
	\index{syarat Mittag-Leffler (Mittag-Leffler condition)}
Misalkan semua morfisme dalam kategori $\mathcal{C}$ mempunyai citra
(Definisi \ref{def:Im-Coim}). Tinjau objek
$X = (X_k, f_k)_{k \geq 0}$ dari $\cate{InvSys}(\mathcal{C})$.
Untuk $b \geq a$, definisikan $f^b_a: X_b \to X_a$ sebagai komposisi
$f_a \cdots f_{b-1}$, dengan konvensi $f^a_a = \identity$. Sistem $X$
dikatakan memenuhi \emph{syarat Mittag-Leffler} apabila
% segment-id: o014.li2.chapter3.example-lim1.q006
\[ \forall k \geq 0, \; \exists N \geq k, \quad n \geq N \implies \Image f^n_k = \Image f^N_k. \]
\end{definisi}

% segment-id: o014.li2.chapter3.example-lim1.p011
Sebagai contoh, jika setiap $f_k$ merupakan epimorfisme, maka
$(X_k, f_k)_{k \geq 0}$ secara otomatis memenuhi syarat Mittag-Leffler.

% segment-id: o014.li2.chapter3.example-lim1.p012
Untuk sistem Mittag-Leffler $(X_k, f_k)_{k \geq 0}$, setiap $X_k$
mempunyai subobjek yang terdefinisi dengan baik
$X_k^\flat := \Image f^N_k$, dengan $N \gg k$. Mudah diperiksa bahwa setiap
morfisme dalam barisan yang bersesuaian
% segment-id: o014.li2.chapter3.example-lim1.q007
\[ \cdots \to X^\flat_2 \xrightarrow{f^\flat_1} X^\flat_1 \xrightarrow{f^\flat_0} X^\flat_0 \]
menjadi epimorfisme, dengan
$f^\flat_k := f_k|_{X^\flat_{k+1}}$.

% segment-id: o014.li2.chapter3.example-lim1.le004
\begin{lema}\label{prop:ML-nonempty}
Ambil $\mathcal{C}$ sebagai $\cate{Set}$ atau $R\dcate{Mod}$, dengan $R$
sebarang gelanggang. Andaikan objek
$(X_n, f_n)_{n \geq 0}$ dari $\cate{InvSys}(\mathcal{C})$ memenuhi
syarat Mittag-Leffler.
% segment-id: o014.li2.chapter3.example-lim1.l001
\begin{enumerate}[(i)]
% segment-id: o014.li2.chapter3.example-lim1.i001
	\item Keluarga morfisme inklusi $X_n^\flat \hookrightarrow X_n$
  menginduksi isomorfisme kanonik
$\varprojlim_n X^\flat_n \rightiso \varprojlim_n X_n$.
% segment-id: o014.li2.chapter3.example-lim1.i002
	\item Dalam kasus $\mathcal{C} = \cate{Set}$, jika setiap $X_n$
  tak kosong, maka $\varprojlim_n X_n \neq \emptyset$.
\end{enumerate}
\end{lema}
% segment-id: o014.li2.chapter3.example-lim1.pf004
\begin{bukti}
Untuk (i), tuliskan konstruksi konkret $\varprojlim$ dalam
kategori-kategori tersebut:
% segment-id: o014.li2.chapter3.example-lim1.q008
\[ \varprojlim_n X_n = \left\{ (x_n)_{n \geq 0} \in \prod_{n \geq 0} X_n : \forall n \geq 0,\; f_n(x_{n+1}) = x_n \right\}; \]
Setiap $x_n$ dalam ruas kanan secara otomatis termasuk dalam
$\bigcap_{N \geq n} \Image(f^N_n)$. Dari sini langsung diperoleh
$\varprojlim_n X_n = \varprojlim_n X^\flat_n$.

Untuk (ii), berdasarkan definisi setiap $X_n^\flat$ tentu tak kosong.
Karena itu, menurut (i), persoalan tereduksi pada kasus ketika setiap
$f_n$ merupakan epimorfisme, dan dalam kasus tersebut pernyataannya
jelas.
\end{bukti}

% segment-id: o014.li2.chapter3.example-lim1.pr001
\begin{proposisi}\label{prop:ML-exactness}
Ambil $R$ suatu gelanggang dan $\mathcal{A} := R\dcate{Mod}$. Jika objek
$A = (A_n, f_n)_{n \geq 0}$ dari
$\cate{InvSys}(\mathcal{A})$ memenuhi syarat Mittag-Leffler, maka
$\lim^1 A = 0$. Secara ekuivalen, untuk setiap barisan eksak pendek
$0 \to A \to B \to C \to 0$ di $\cate{InvSys}(\mathcal{A})$, barisan
yang bersesuaian $0 \to \varprojlim A \to \varprojlim B \to
\varprojlim C \to 0$ juga eksak.
\end{proposisi}
% segment-id: o014.li2.chapter3.example-lim1.pf005
\begin{bukti}
Pertama, buktikan ekuivalensi kedua pernyataan tersebut. Misalkan
$A \in \cate{InvSys}(\mathcal{A})$. Jika $\lim^1 A = 0$, keeksakan
$0 \to \varprojlim A \to \varprojlim B \to \varprojlim C \to 0$
merupakan akibat langsung dari barisan eksak panjang. Sebaliknya,
andaikan $0 \to \varprojlim A \to \varprojlim B \to \varprojlim C \to 0$
selalu eksak. Pilih monomorfisme $A \hookrightarrow B$ sedemikian
sehingga $B$ merupakan objek injektif dari
$\cate{InvSys}(\mathcal{A})$. Pergeseran dimensi (Proposisi
\ref{prop:dimension-shifting}) kemudian memberikan
$\mathrm{R}^1 \varprojlim A = 0$, dan Teorema \ref{prop:lim1}
selanjutnya menghasilkan $\lim^1 A = 0$.

Sekarang buktikan pernyataan utama. Misalkan $A$ dalam barisan eksak
pendek
$0 \to \underbracket{(A_n, f_n)_n}_{= A} \xrightarrow{\varphi} \underbracket{(B_n, g_n)_n}_{= B} \xrightarrow{\psi} \underbracket{(C_n, h_n)_n}_{= C} \to 0$
memenuhi syarat Mittag-Leffler.
\footnote{Catatan penerjemah (O014-C063): pada sumber, pasangan sistem $C$ di bawah garis kurung ditulis sebagai $(C_n,h_n)$ tanpa subskrip indeks. Karena $C$ merupakan objek dari $\cate{InvSys}(\mathcal{A})$, target memulihkan notasi tersebut menjadi $(C_n,h_n)_n$.}
Tujuan kita ialah membuktikan bahwa $\varprojlim \psi$ merupakan
epimorfisme. Berdasarkan Lema \ref{prop:ML-nonempty} (ii), cukup dibuktikan
bahwa, untuk setiap $c = (c_n)_{n \geq 0} \in \varprojlim C$, sistem invers
himpunan tak kosong
% segment-id: o014.li2.chapter3.example-lim1.q009
\[ \cdots \to \psi_2^{-1}(c_2) \xrightarrow{g_1} \psi_1^{-1}(c_1) \xrightarrow{g_0} \psi_0^{-1}(c_0) \]
memenuhi syarat Mittag-Leffler; hal ini akan menunjukkan bahwa
$(\varprojlim \psi)^{-1}(c)$ tak kosong.

Diberikan $k \geq 0$, pilih $N \geq k$ yang cukup besar sehingga
$\Image f^n_k = \Image f^N_k$ untuk $n \geq N$. Pilih
$b_N \in \psi_N^{-1}(c_N)$. Untuk memeriksa syarat Mittag-Leffler,
kita akan menunjukkan bahwa, untuk setiap $n \geq N$, terdapat
$b_n \in \psi_n^{-1}(c_n)$ sedemikian sehingga
$g^n_k(b_n) = g^N_k(b_N)$.

Pilih sebarang $b'_n \in \psi_n^{-1}(c_n)$. Karena
$\psi_N\left( g^n_N(b'_n)\right) = c_N =
\psi_N\left(b_N \right)$, terdapat $a_N$ sedemikian sehingga
$g^n_N(b'_n) - b_N = \varphi_N(a_N)$. Selanjutnya, pilih $a_n$
sedemikian sehingga $f^n_k(a_n) = f^N_k(a_N)$. Maka,
$b_n := b'_n - \varphi_n(a_n)$ mempunyai sifat yang diinginkan.
\end{bukti}

% segment-id: o014.li2.chapter3.example-lim1.p013
Untuk kategori abelian umum $\mathcal{A}$ yang memenuhi Hipotesis
\ref{hyp:lim1}, terdapat contoh tandingan yang menunjukkan bahwa syarat
Mittag-Leffler tidak selalu mengakibatkan $\lim^1 = 0$; lihat
\cite{Roo06} untuk pembahasan terperinci.

% segment-id: o014.li2.chapter3.example-lim1.p014
Hasil berikut merupakan latihan sederhana yang akan digunakan dalam
\S\sourcecrossref{sec:K-injectives}{3.15}.

% segment-id: o014.li2.chapter3.example-lim1.co001
\begin{korolari}\label{prop:ML-4-terms}
Misalkan $\mathcal{A}$ memenuhi Hipotesis \ref{hyp:lim1}, dan misalkan
$A \to B \to C \to D$ merupakan barisan eksak di
$\cate{InvSys}(\mathcal{A})$. Jika $\lim^1 A = 0$, maka
% segment-id: o014.li2.chapter3.example-lim1.q010
\[ \varprojlim B \to \varprojlim C \to \varprojlim D \]
merupakan barisan eksak di $\mathcal{A}$.
\end{korolari}
% segment-id: o014.li2.chapter3.example-lim1.pf006
\begin{bukti}
Tetapkan $H := \Ker[C \to D]$ dan $X := \Image[A \to B]$. Karena
$\varprojlim$ eksak kiri, terdapat isomorfisme kanonik
% segment-id: o014.li2.chapter3.example-lim1.q011
\[ \Ker\left[ \varprojlim C \to \varprojlim D\right] \simeq \varprojlim H . \]
Morfisme $B \to C$ terfaktorkan sebagai
$B \twoheadrightarrow H \hookrightarrow C$. Karena itu, persoalan
tereduksi pada pembuktian bahwa
$\varprojlim B \to \varprojlim H$ merupakan epimorfisme. Contoh
\ref{eg:quotient-lim1-0} menunjukkan bahwa $\lim^1 X = 0$. Dengan
menerapkan $\varprojlim$ pada barisan eksak pendek
$0 \to X \to B \to H \to 0$, diperoleh bahwa
$\varprojlim B \to \varprojlim H$ merupakan epimorfisme.
\end{bukti}
